\section{Finite-chain separations}\label{sec:finite-separation}

Three constructions distinguish finite critical-chain resources from
every maximally entangled state. The original game has an assisted
certificate at $N=2^{19}$, using the chain together with one shared
nine-dimensional maximally entangled register, and a chain-only
certificate at $N=2^{84}$. A four-spin construction uses a tagged
tilted CHSH test and the chain ground state alone.

\subsection{The original game with a shared nine-dimensional register}
\label{sec:xx-assisted-separation}

\begin{theorem}[Assisted finite XX separation]\label{thm:xx-assisted-separation}
Let $\psi_N$ be the ground state of the half-filled open XX chain on
$N$ sites, cut in the middle. In the original uniform game $G$,
$\psi_{2^{19}}\otimes\Phi_9$ admits a strategy with loss below
$0.000703$, and $\psi_{2^{20}}\otimes\Phi_9$ admits one with loss below
$0.000663$. Both losses are strictly below
$\gamma=(3-2\sqrt2)/243$, the floor of Theorem~\ref{thm:operational-tracial}
for every maximally entangled state of arbitrary finite rank and arbitrary
local POVMs.
\end{theorem}

\begin{proof}
The strategy of Theorem~\ref{thm:operational-upper} on
$\psi_N\otimes\Phi_9$ has loss exactly $2\cE(\psi_N)/243$.
Appendix~\ref{app:xx-assisted-certificate} bounds this expression using
the exact finite XX covariance, certified quadrature and spectral errors,
and directed integer rank calculations. The resulting upper bounds are
$0.000702783850$ and $0.000662385080$, respectively.
\end{proof}

The auxiliary register is part of the shared input. Providing the same
register to a maximally entangled competitor gives
$\Phi_D\otimes\Phi_9\simeq\Phi_{9D}$.
The prescribed assisted strategy beats $\gamma$ exactly when
\[
 \cE(\psi_N)<\frac{3-2\sqrt2}{2}=0.0857864376\ldots.
\]
The numerical estimates for $\cE(\psi_N)$ at $N=2^{18},2^{19},2^{20}$
round to $0.0908$, $0.0854$, and $0.0805$, respectively. The certified
intervals in Table~\ref{tab:xx-assisted-certificate} lie above the threshold
at $N=2^{18}$ and below it at the other two sizes.

\subsection{An explicit finite chain for the original game}\label{sec:xx84-separation}

\begin{theorem}\label{thm:xx84-separation}
The ground state of the open XX chain of length $N=2^{84}$, cut in the
middle, has a strategy in the uniform game $\GIII$ with loss below
$0.000461$. The chain state is the strategy's sole shared
resource. This loss is strictly below the floor for every maximally
entangled state in Theorem~\ref{thm:operational-tracial}.
\end{theorem}

The chain length is far beyond the sizes of the numerical sweep. The
certificate uses a $1418\times1418$ approximation to one parity block
of the exact cross covariance and integer rank counts for a finite model
of its Schmidt spectrum. The dimension of the covariance calculation
grows with the number of dyadic intervals in the half-chain. The model
strategy has loss below $0.000410720762$. Propagating the spectral
approximation and arithmetic bounds gives the certified loss bound in
the theorem for the physical ground state. Appendix~\ref{app:xx84-certificate}
gives the complete calculation.

\subsection{A four-spin separation with tilted CHSH}
\label{sec:finite-chain-hybrid}

The four-spin construction realizes the known tilted CHSH separation on
an XX-chain ground state~\cite{AMP12,BP15}. For $0<\alpha\le2$, define
$T_\alpha$ with binary answers $a,b\in\{-1,1\}$.
Each of the four pairs $(x,y)\in\{0,1\}^2$ is asked with probability
$1/(4+\alpha)$ and wins when $ab=(-1)^{xy}$. The pair $(0,*)$ is asked
with probability $\alpha/(4+\alpha)$ and wins when $a=1$. The question
$0$ is the same in both rules. If $A_x,B_y$ denote the differences of the
two answer effects, the loss is $(4+\alpha-S_\alpha)/(2(4+\alpha))$ where
\begin{equation}\label{eq:tilted-score}
 S_\alpha=\left\langle\alpha A_0+A_0(B_0+B_1)+A_1(B_0-B_1)\right\rangle.
\end{equation}
Let $H_\alpha$ choose $\GIII$ with probability $1/10$ and $T_\alpha$ with
probability $9/10$. Each player receives the branch tag as part of the question.

\begin{theorem}\label{thm:finite-chain-hybrid}
The ground state of the four-site open XX chain, cut after site two, has a
strategy on $H_1$ with loss less than $0.177$. Every maximally entangled state,
of any finite rank and with arbitrary local POVMs, has loss at least $0.180$.
The ground state is the sole shared resource, and the strategy plays $G$
classically. The same loss floor holds for mixtures of maximally entangled
states.
\end{theorem}

\begin{proof}
On $\Phi_d$, absorb Bob's transpose into his effects and write
$\tau=\Tr/d$. For fixed $B_0,B_1$, maximizing over the selfadjoint
contractions $A_0,A_1$ allows them to be reflections. Indeed,
$\max_{-I\le A\le I}\tau(AK)=\tau|K|$ for every selfadjoint $K$, and the
maximum is attained by the spectral sign of $K$ with either sign on its kernel.
Optimizing over Bob's two contractions now yields
\[
 \alpha\tau(A_0)+\tau|A_0+A_1|+\tau|A_0-A_1|.
\]
The selfadjoint operator $K=(A_0A_1+A_1A_0)/2$ commutes with the reflections.
On its $k$ eigenspace, $|A_0\pm A_1|=\sqrt{2\pm2k}\,I$.
For $-1<k<1$, the reflection $(A_1-kA_0)/\sqrt{1-k^2}$ anticommutes
with $A_0$, so $\Tr A_0=0$ on this eigenspace. Its contribution per
dimension is at most $2\sqrt2$. At $k=\pm1$ the contribution is at most
$2+\alpha$. Thus
\begin{equation}\label{eq:tilted-tracial}
 S_\alpha(\Phi_d)\le\max\{2\sqrt2,2+\alpha\}.
\end{equation}
Local effects can always be compressed to the Schmidt supports, so this
argument covers arbitrary local POVMs. Linearity for fixed measurements
extends the corresponding loss bound to mixtures of flat states.

Use the hopping convention $-\sum_{j=1}^3(a_j^*a_{j+1}+a_{j+1}^*a_j)$ for
the open XX chain after Jordan--Wigner transformation. Reversing the hopping
sign is an on-site unitary change of basis. The occupied standing waves are
$v_1=\sqrt{2/5}\,(s,t,t,s)$ and $v_2=\sqrt{2/5}\,(t,s,-s,-t)$, where
$s=\sin(\pi/5)$ and $t=\sin(2\pi/5)$. In the left half their covariance is
\[
 \frac25\begin{pmatrix}s^2+t^2&2st\\2st&s^2+t^2\end{pmatrix}
   =\begin{pmatrix}1/2&1/\sqrt5\\1/\sqrt5&1/2\end{pmatrix}.
\]
In particular $\nu=1/2+1/\sqrt5$ is its larger eigenvalue. Define the
left orbitals $\ell_1=(1,1)/\sqrt2$, $\ell_2=(-1,1)/\sqrt2$ and the
right orbitals $r_1=(1,-1)/\sqrt2$, $r_2=(1,1)/\sqrt2$. The occupied
vectors $(v_1+v_2)/\sqrt2$ and $(v_1-v_2)/\sqrt2$ are respectively
$\sqrt\nu\,\ell_1+\sqrt{1-\nu}\,r_1$ and
$\sqrt{1-\nu}\,\ell_2+\sqrt\nu\,r_2$.

Use local occupation bases in these orbitals with left creation operators
ordered before right ones. If $(n_1,n_2)$ are the left occupations in the
two occupied orbital pairs, the right occupations are $(1-n_1,1-n_2)$
and the Slater determinant has phase $(-1)^{(1-n_1)n_2}$ in that component,
up to an overall phase. The local basis changes
\[
 \ket{n_1n_2}_A\longmapsto (-1)^{(1-n_1)n_2}\ket{1-n_1,n_2}_A,
 \qquad
 \ket{m_1m_2}_B\longmapsto\ket{m_1,1-m_2}_B
\]
therefore give exactly two Schmidt qubit pairs, both in
$\ket{\xi}=\sqrt\nu\ket{00}+\sqrt{1-\nu}\ket{11}$.
On the first pair measure $A_0=Z$, $A_1=X$ and
$B_0=(\sqrt5Z+X)/\sqrt6$, $B_1=(\sqrt5Z-X)/\sqrt6$. Its three
expectations $\langle Z\otimes I\rangle$, $\langle Z\otimes Z\rangle$,
and $\langle X\otimes X\rangle$ are $2/\sqrt5$, $1$, and $1/\sqrt5$.
The score in~\eqref{eq:tilted-score} is consequently
\begin{equation}\label{eq:tilted-xx4-score}
 S_\alpha(\psi_4)\ge \frac{2\sqrt6+2\alpha}{\sqrt5}.
\end{equation}
The other pair is ignored. In the $\GIII$ branch use the deterministic
strategy of Theorem~\ref{thm:operational-upper}. With $\alpha=1$ the
total loss on $\psi_4$ is at most
\[
 \frac1{270}+\frac9{100}\left(5-\frac{2\sqrt6+2}{\sqrt5}\right)
 <\frac{177}{1000}.
\]
For example, $\sqrt6>2449/1000$ and $\sqrt5<2237/1000$ show
$(2\sqrt6+2)/\sqrt5>37/12$. Equation~\eqref{eq:tilted-tracial}
gives loss at least $(9/100)(5-3)=9/50$ on every $\Phi_d$.
\end{proof}

\paragraph{Question weights.}
The separation also holds with the rational weight $\alpha=5/6$.
One can draw the four CHSH pairs with six units of weight each and the
marginal pair with five units. The tracial score bound is $17/6$ and the
four-site strategy gives score at least $(2\sqrt6+5/3)/\sqrt5$.
The total loss is less than $0.181$ on the ground state and at least
$27/145>0.186$ on all flat states. For the fixed four-site observables in
the proof, the unnormalized score gap over~\eqref{eq:tilted-tracial} is
largest at $\alpha=2\sqrt2-2$.

\paragraph{Stability under preparation error.}
For either question distribution, write $L_{\rm flat}$ for the floor from
\eqref{eq:tilted-tracial} and $L_4$ for the achieved four-site loss. The
exact formulas give $L_{\rm flat}-L_4>0.0039$ at $\alpha=1$ and
$L_{\rm flat}-L_4>0.0058$ at $\alpha=5/6$. Both inequalities follow by
using $\sqrt6>2449489/10^6$ and $\sqrt5<2236068/10^6$ in the score.
Trace distance $\eta$ from the four-site state raises the loss of its
fixed strategy by at most $\eta$. Thus the separation survives preparation
error in trace distance up to $0.0039$ or $0.0058$, respectively.

\paragraph{The optimal benchmark and the type obstruction.}
Ac\'in, Massar and Pironio give the sharp tilted CHSH score
$\sqrt{8+2\alpha^2}$~\cite{AMP12}, and Bamps and Pironio give
sum-of-squares decompositions~\cite{BP15}. For commuting local reflections
$A_0,A_1,B_0,B_1$, set $C=(B_0+B_1)/2$, $D=(B_0-B_1)/2$,
$U=(A_0A_1+A_1A_0)/2$ and $V=(A_0A_1-A_1A_0)/(2i)$.
Then $C^2+D^2=I$, $CD+DC=0$, and $U^2+V^2=I$. If $W$ is the operator
in~\eqref{eq:tilted-score},
\[
 W^2=(4+\alpha^2)I+R,\qquad R=4\alpha C+4\alpha U D+8V(iCD).
\]
The three terms of $R$ anticommute pairwise. For the commuting
positive contractions $u=U^2$ and $c=C^2$ this gives
\[
 \frac{R^2}{16}=\alpha^2u+(1-u)\bigl[\alpha^2c+4c(1-c)\bigr]
 \le \frac{(4+\alpha^2)^2}{16}I.
\]
The scalar bound follows by maximizing the quadratic on $0\le c\le1$
and using $\alpha^2\le(4+\alpha^2)^2/16$. Hence $\norm W^2\le8+2\alpha^2$.
Every selfadjoint contraction is an average of its spectral threshold
reflections over thresholds in $[-1,1]$. Since $W$ is affine in each
observable, the same score bound holds for contractions in commuting
local von Neumann algebras. Equality for $0<\alpha<2$ is achieved on
$\cos\theta\ket{00}+\sin\theta\ket{11}$ with
$\sin(2\theta)=\sqrt{(4-\alpha^2)/(4+\alpha^2)}$, using $A_0=Z$, $A_1=X$
and $B_{0,1}=(Z\pm\sin(2\theta)X)/\sqrt{1+\sin^2(2\theta)}$.
The endpoint $\alpha=2$ is attained on a product state.

Consequently the loss of $H_\alpha$ is at least
\[
 L_\alpha^* = \frac9{10}\frac{4+\alpha-\sqrt{8+2\alpha^2}}{2(4+\alpha)}.
\]
Equality forces a perfect strategy in the $\GIII$ branch. In particular,
the original obstruction to exact play in the presence of a semifinite
local trace applies to attaining this benchmark.
For $0<\alpha<2$, the ground states of the finite open XX chains approach
$L_\alpha^*$ as $N\to\infty$. The $\GIII$ branch approaches perfect play
by Theorem~\ref{thm:value}. For $T_\alpha$, the exact finite-chain
one-body entanglement energies contain a mode tending to the strictly
positive energy $\log\cot^2\theta$. Indeed, in any fixed neighbourhood
of this energy there are order $\log N$ mesh entries from
Corollary~\ref{cor:xx-mesh}. The bounded sum of singular-value errors
forces at least one true mode in the neighbourhood for all sufficiently
large $N$. Measuring that mode approaches the optimal two-qubit score.
The two branch strategies can be selected independently by the question tags.
